\subsubsection{AI এজেন্ট: লক্ষ্য থেকে কাজ}
\definition{AI এজেন্ট}{লক্ষ্য অর্জনের জন্য পরিবেশের তথ্য নিয়ে সিদ্ধান্ত ও কাজ করে এমন ব্যবস্থাকে AI Agent বলে। আধুনিক LLM-ভিত্তিক এজেন্টে মডেল কাজের ধাপ ও টুল নির্বাচন করতে পারে; সফটওয়্যার অনুমোদিত টুল চালিয়ে ফল ফেরত দেয়।}
এজেন্টের ধারণা কথোপকথনভিত্তিক AI-এর আগেও ছিল; রোবট ও গেমের এজেন্ট উদাহরণ। এটি একটি ব্যবস্থার কাজ করার ধরন, শুধু নতুন মডেল পরিবারের নাম নয়। টুল চালানোর নির্দেশ লেখা এবং সফটওয়্যারের সেই কাজ করা পৃথক ধাপ।
\chOnePair{\chOneNative{AgentDetail}}{লক্ষ্য ও সীমা অনুযায়ী পরের কাজ নির্বাচন করা হয়। টুল বা পরিবেশে কাজের ফল পর্যবেক্ষণ হিসেবে ফিরে আসে। সেই ফল ও আগের অগ্রগতি নিয়ে পরবর্তী ধাপ নির্ধারিত হয়। লক্ষ্য পূরণ, সীমা অতিক্রম বা মানুষের সাহায্যের প্রয়োজন হলে চক্র থামে।}
\begin{compactenum}
\item লক্ষ্য, অনুমোদিত ফাইল ও সেবা এবং সময় ও ব্যয়ের সীমা নির্ধারণ করা হয়।
\item মডেল বা নিয়ন্ত্রক প্রাপ্ত তথ্য থেকে কাজ নির্বাচন করে; প্রয়োজনে কাজকে ছোট ধাপে ভাগ করে। সব এজেন্টে দীর্ঘ লিখিত পরিকল্পনা থাকে না।
\item অনুমোদিত অনুসন্ধান, ক্যালকুলেটর, ফাইল-পাঠক, তথ্যভান্ডার বা কোড চালানোর টুল ব্যবহার করা হয়।
\item কাজের অগ্রগতি, তথ্য ও টুলের ফল রাখা হয়। এই স্মৃতি সরাসরি স্থায়ী প্রশিক্ষণ বদলায় না।
\item সফলতা বা ব্যর্থতার ফল দেখে ধাপ সংশোধন করা হয়।
\item প্রত্যাশিত ফল, ভুলের সীমা, মানুষের পর্যালোচনা ও থামার শর্ত মিলিয়ে কাজ শেষ করা হয়।
\end{compactenum}
সাধারণ চ্যাটবট প্রশ্ন থেকে উত্তর দেয়; টুল থাকা ঐচ্ছিক। নির্দিষ্ট Workflow-এ আগে থেকে লেখা ধাপ চলে। LLM এজেন্ট ফল দেখে পরের ধাপ বা টুল নির্বাচন করতে পারে। শারীরিক রোবটের এজেন্ট সেন্সরের তথ্য নিয়ে বাস্তব পরিবেশে কাজ করে। শুধু টুল আছে বলেই পূর্ণ স্বায়ত্তশাসন প্রমাণিত হয় না; সব কাজ সফল বা স্বাধীনও হতে হবে না।
\subtopichead{কলেজের ফলাফলের প্রতিবেদন}
শিক্ষক বিভিন্ন শ্রেণির CSV থেকে মোট, গড় ও অনুপস্থিত ফলের তালিকা চান। এজেন্ট অনুমোদিত ফাইল পড়ে কলাম শনাক্ত করে এবং হিসাবের টুল দিয়ে ফল বের করে। অনুপস্থিত মান আলাদা রেখে পরিচয় বাদ দিয়ে খসড়া তৈরি করা যায়। শিক্ষক নমুনা হিসাব মিলিয়ে দেখেন। ভুল কলাম বা অনুপস্থিত ফাইল পেলে অনুমান করে পূর্ণ ফল বানানো ঠিক নয়; হিসাবের টুল ও ভাষা মডেলের ভূমিকা আলাদা।
\subtopichead{কোড সংশোধনের উদাহরণ}
HTML পাতায় ছবি না এলে এজেন্ট অনুমোদিত ফাইল ও ছবির পথ দেখতে, সংশোধন করতে এবং ব্রাউজার বা পরীক্ষার ফল পড়তে পারে। ভুল থাকলে আবার পরীক্ষা করে। কেবল কোডের পরামর্শ চ্যাটবটও দিতে পারে; ফাইল দেখা, সংশোধন, পরীক্ষা এবং ফল অনুযায়ী পরের কাজ নির্বাচন এজেন্টের কার্যক্রম বোঝায়।
\subtopichead{সুবিধা, সীমা ও নিরাপত্তা}
এজেন্ট বিভিন্ন টুলের তথ্য একত্র করে বহু ধাপের পুনরাবৃত্ত কাজ করতে সাহায্য করে। ভুল সিদ্ধান্ত, ভুল টুলের নির্দেশ বা অবিশ্বস্ত তথ্যের প্রভাব পরের ধাপে ছড়াতে পারে। অপ্রয়োজনীয় চক্র সময় ও ব্যয় বাড়ায়। \term{Multi-agent System}-এ একাধিক এজেন্ট কাজ ভাগ করে বা পরস্পরের ফল পর্যালোচনা করে; সমন্বয় ও গণনার ব্যয় বাড়ে এবং একাধিক এজেন্টও সঠিক ফল নিশ্চিত করে না।
\begin{compactpoints}
\item কাজের প্রয়োজন অনুযায়ী সীমিত ফাইল, তথ্য ও টুলের অধিকার দিতে হয়।
\item বাইরের ওয়েবপাতা বা নথির নির্দেশ ব্যবহারকারীর অনুমতি নয়। অবিশ্বস্ত লেখায় আচরণ বদলানোর চেষ্টা \term{Prompt Injection}-এর উদাহরণ।
\item প্রকাশ, টাকা পাঠানো বা রেকর্ড বদলানোয় অনুমোদন ও পর্যালোচনার নিয়ম দরকার। খসড়া তৈরি ও প্রকাশ পৃথক কাজ।
\item টুলের লগ, পরীক্ষার ফল ও অসম্পূর্ণ কাজের অবস্থা রাখতে হয়। সুন্দর প্রতিবেদনের বিন্যাস দিয়ে ভুল হিসাব সঠিক হয় না।
\end{compactpoints}
\subsubsection{অবদানের সীমা ও ভবিষ্যৎ উন্নয়ন}
AI-এর সুফল কাজ, প্রতিষ্ঠান, দেশ ও ব্যবহারের ধরন অনুযায়ী বদলায়। একটি গ্রাহকসেবা গবেষণায় সহায়ক AI ব্যবহারে গড় উৎপাদনশীলতা বেড়েছিল; কম অভিজ্ঞ কর্মীর সুফল বেশি ছিল। একটি প্রতিষ্ঠানের ফল সব কর্মী বা খাতের একই উন্নতির হার বোঝায় না।
AlphaFold প্রোটিনের গঠন সম্পর্কে পূর্বাভাস দিয়ে গবেষকের সম্ভাব্য প্রশ্ন ও পরীক্ষা বাছাইয়ে সাহায্য করে। এটি AI ও বায়োইনফরমেটিকসের সংযোগের উদাহরণ। মডেলের পূর্বাভাস এবং পরীক্ষাগারের প্রমাণ পৃথক; গঠন সম্পর্কে ফল পাওয়া মানেই ওষুধের কার্যকারিতা বা চিকিৎসা প্রমাণিত হওয়া নয়।
ভবিষ্যৎ উন্নয়ন ভালো ডেটা, গণনাক্ষমতা, পরীক্ষা, ভাষাসমর্থন, নিরাপত্তা, ব্যয় ও মানুষের গ্রহণের ওপর নির্ভর করবে। নিচের দিকগুলো সম্ভাবনা, নির্দিষ্ট সময়ের নিশ্চয়তা নয়।
\begin{compactpoints}
\item বহুমাধ্যম AI লেখা, কথা, ছবি ও ভিডিওর সম্পর্ক ব্যবহার করে শেখা ও সহায়ক প্রযুক্তি উন্নত করতে পারে। প্রতিটি মডেলে সব ইনপুট ও আউটপুট থাকবে এমন নিশ্চয়তা নেই।
\item ছোট ও স্থানীয় মডেল মোবাইল বা সীমিত সংযোগে কিছু কাজ সহজ করতে পারে। গতি, মেমরি, বিদ্যুৎ ও নির্ভুলতার সমন্বয় দরকার।
\item বাংলার ভালো ডেটা, বিভিন্ন উচ্চারণ ও স্থানীয় কাজের মূল্যায়ন দরকার। বিদেশি উদাহরণে ভালো ফল বাংলাদেশের সব কাজের উপযোগিতা প্রমাণ করে না।
\item সক্ষম এজেন্ট ও রোবটের কাজ বাড়তে পারে; দীর্ঘ কাজের নির্ভরযোগ্যতা, কাজের সীমা ও পরিবেশে পরীক্ষা উন্নত করতে হবে।
\item গবেষণার পথ বাছাই ও কিছু পুনরাবৃত্ত কাজ দ্রুত হতে পারে। যাচাই, ডেটা ও ব্যবস্থা পরিচালনার দক্ষতা লাগবে; চাকরির পরিবর্তনের নিশ্চিত সংখ্যা বলা যায় না।
\item AGI বা Artificial General Intelligence-এর সংজ্ঞা ও মূল্যায়নে বিতর্ক আছে। একটি পরীক্ষায় ভালো ফল বা সাবলীল কথোপকথন দিয়ে বিস্তৃত সক্ষমতা বা চেতনা নিশ্চিত হয় না; নির্দিষ্ট বছরের নিশ্চয়তাও নেই।
\end{compactpoints}
শিক্ষার্থীর বিষয়জ্ঞান, স্পষ্ট নির্দেশ, ভালো ডেটা চেনা, মৌলিক প্রোগ্রামিং ও হিসাব এবং ফল পরীক্ষা করার দক্ষতা দরকার। খসড়া ব্যবহার করলেও নিজের ভাষায় অর্থ ব্যাখ্যা এবং প্রমাণ দিয়ে ভুল সংশোধন করতে হবে।
